\chapter[Harmonic Admission Scheduling]{Harmonic Admission Scheduling for Computational Experiments}
\label{app:harmonic-admission}

This appendix records an engineering control used while running the computational parts of the research program. Its evidential class is \texttt{ENGINEERING\_CONTROL}: it regulates when local processes may start so that numerical work, rendering, repository I/O, and document compilation are less likely to collide in the same execution window. The scheduling convention does not participate in the derivation of the temporal field equations.

\section{Dyadic admission bands}

A base cadence
\[
f_0 = 7.83\;\mathrm{Hz},
\qquad
T_0=f_0^{-1},
\]
is used as an implementation clock. The workload bands are defined by
\[
f_k=\frac{f_0}{2^k},
\qquad
T_k=\frac{2^k}{f_0}.
\]
Band $B_0$ is reserved for lightweight heartbeat observation. Workload admission begins at $B_1$.

For integer base tick $n$, the phase class of band $B_k$, $k\geq1$, is
\[
n\equiv 2^{k-1}\pmod{2^k}.
\]
These classes are pairwise disjoint. Consequently, a base tick can admit at most one workload class among $B_1,\ldots,B_K$.

\begin{center}
\begin{tabular}{ccl}
\hline
Band & Frequency & Operational role \\
\hline
$B_0$ & $7.830$ Hz & heartbeat observation \\
$B_1$ & $3.915$ Hz & queue and admission \\
$B_2$ & $1.9575$ Hz & candidate microsteps \\
$B_3$ & $0.97875$ Hz & critic and graph audit \\
$B_4$ & $0.489375$ Hz & numerics and small slices \\
$B_5$ & $0.2446875$ Hz & serialization and repository I/O \\
$B_6$ & $0.12234375$ Hz & visualization \\
$B_7$ & $0.061171875$ Hz & LaTeX and heavy builds \\
\hline
\end{tabular}
\end{center}

The numerical value of $f_0$ is an operational scheduling convention in this appendix. Any physical use of a frequency scale requires an independent claim, derivation, and evidence path in the main formalism.

\section{Skip semantics and jitter control}

An admission window is permission to start queued work rather than a demand for periodic execution. Empty windows perform no work. A missed window is skipped:
\[
\text{missed slot}\longrightarrow\text{skip},
\]
so delayed processes cannot produce a catch-up burst. If a task assigned to $B_k$ is still active at its next eligible slot, the new instance is skipped as well.

Measured scheduling jitter $J_k$ can demote a workload to a slower band,
\[
J_k>J_{\max}
\quad\Longrightarrow\quad
B_k\mapsto B_{k+1},
\]
subject to the configured maximum band. This mechanism is intended to protect the interactive research workflow from render, I/O, and compilation spikes while keeping the computational procedures reproducible.

\section{Reference implementation}

The reference implementation is stored in \texttt{src/idt/harmonic\_scheduler.py}. The scheduler exposes the band definitions, exact phase-slot predicate, next-slot calculation, quiet deep-boundary slots, and jitter-triggered demotion. Reference tests verify pairwise disjoint workload classes and the no-catch-up admission semantics.
